% --- Plik: thesis-packages.tex ---

%% Do kompilacji używaj XeLaTeX!
\usepackage{fontspec}

% Domyślna czcionka dokumentu (zamiennik dla Helvetica wg wytycznych)
\setmainfont{Arial}
\setsansfont{Arial}
\renewcommand{\familydefault}{\sfdefault}

% Definiowanie niestandardowych czcionek uczelnianych
% (Nazwy w nawiasach klamrowych muszą odpowiadać tym zainstalowanym w Twoim systemie)
\newfontfamily\adagioslab{Adagio_Slab-Regular}
\newfontfamily\adagioslablight{Adagio_Slab-Light}

\usepackage{polyglossia}
\setdefaultlanguage{polish}
\setotherlanguage{english}

% Ustawienia marginesów: wewnętrzny 30 mm, zewnętrzny 20 mm, górny i dolny 25 mm [3]
\usepackage[inner=30mm, outer=20mm, top=25mm, bottom=25mm]{geometry}

% Interlinia 1.15 [3]
\usepackage{setspace}
\setstretch{1.15}

% Pakiety dodatkowe
\usepackage{graphicx}
\usepackage{fancyhdr}
\usepackage{url}
\usepackage{xcolor}    % Pakiet do obsługi kolorów
\usepackage{hyperref}  % Pakiet do obsługi hiperłączy

\hypersetup{
	colorlinks=true,   % Włącza kolorowanie tekstu zamiast rysowania ramek
	linkcolor=black,   % Odnośniki wewnętrzne (np. spis treści) pozostają czarne do druku
	citecolor=black,   % Odnośniki do bibliografii pozostają czarne
	urlcolor=blue      % Adresy URL i e-mail będą wyraźnie niebieskie
	%urlcolor=blue!80!black
}


\usepackage{titlesec}
\usepackage{caption} % Do formatowania podpisów
\usepackage{booktabs}
\usepackage[table]{xcolor} % Włącza zaawansowane kolorowanie tabel

% Pakiety do obsługi zaawansowanej matematyki
\usepackage{amsmath}
\usepackage{amsfonts}


\usepackage{listings} % Do obsługi kodów źródłowych programów (np. Python, C++, MATLAB)
% Kolory do kodów Pythona w Jupyter Notebook
\definecolor{ju_keyword}{rgb}{0.0, 0.5, 0.0}
\definecolor{ju_string}{rgb}{0.73, 0.13, 0.13}
\definecolor{ju_comment}{rgb}{0.25, 0.50, 0.50}
\definecolor{ju_builtin}{rgb}{0.0, 0.30, 0.60}
\definecolor{ju_func}{rgb}{0.0, 0.0, 1.0}
\definecolor{ju_tlo}{rgb}{0.99, 0.99, 0.99}
\definecolor{ju_linia}{rgb}{0.85, 0.85, 0.85}
\lstset{
	language=Python,
	alsoletter={_},
	backgroundcolor=\color{ju_tlo},
	basicstyle=\ttfamily\small\color{black},
	keywordstyle=\bfseries\color{ju_keyword},
	commentstyle=\itshape\color{ju_comment},
	stringstyle=\color{ju_string},
	morekeywords=[2]{print,len,range,str,int,float,list,dict,set,tuple,zip,enumerate,abs,min,max,sum},
	keywordstyle=[2]\color{ju_builtin},
	emph={metoda_trapezow,metoda_simpsona,wykres_bledow,moja_funkcja},
	emphstyle=\color{ju_func},
	frame=single,
	rulecolor=\color{ju_linia},
	numbers=left,
	numberstyle=\tiny\color{gray!70},
	breaklines=true,
	columns=flexible,
	keepspaces=true,
	xleftmargin=15pt,
	xrightmargin=5pt,
	literate={ą}{{\k{a}}}1 {ć}{{\'c}}1 {ę}{{\k{e}}}1 {ł}{{\l}}1 {ń}{{\'n}}1 {ó}{{\'o}}1 {ś}{{\'s}}1 {ź}{{\'z}}1 {ż}{{\.z}}1
}

% Ciągła numeracja wzorów przez całą pracę
\counterwithout{equation}{chapter}

% Ciągła numeracja tabel przez całą pracę
\counterwithout{table}{chapter}

% Ciągła numeracja rysunków przez całą pracę
\counterwithout{figure}{chapter}

\AtBeginDocument{
    \counterwithout{lstlisting}{chapter} % Ciągła numeracja listingów przez całą pracę
	\renewcommand{\lstlistingname}{Kod źródłowy}   % Zmienia "Listing 1" na "Kod źródłowy 1"
}


% Konfiguracja stylu stron, nagłówków i stopek (Pakiet fancyhdr)
% Numeracja stron: na dole, po zewnętrznej stronie, druk dwustronny
\pagestyle{fancy}
\fancyhf{}
\fancyfoot[LE,RO]{\thepage}        % Bieżący numer strony: lewa parzysta (LE), prawa nieparzysta (RO)
\renewcommand{\headrulewidth}{0pt} % Całkowite usunięcie poziomej linii z nagłówka

% Zapewnienie spójności formatowania na stronach otwierających rozdziały 
% Domyślnie LaTeX wymusza na początku każdego rozdziału styl 'plain' (numer na środku).
% Poniższy blok nadpisuje styl 'plain', przenosząc numer na krawędzie boczne (LE/RO).
\fancypagestyle{plain}{%
	\fancyhf{}%
	\fancyfoot[LE,RO]{\thepage}%
	\renewcommand{\headrulewidth}{0pt}%
}

% Rozwiązanie konfliktu z pakietem formatowania tytułów (titlesec)
% Pakiet 'titlesec' często nadpisuje zachowanie komendy \chapter i ignoruje styl 'plain'.
% Poniższa instrukcja bezwzględnie wymusza stosowanie stylu 'fancy' na stronach z \chapter.
\usepackage{titlesec}
\assignpagestyle{\chapter}{fancy}


% Formatowanie nagłówków rozdziałów (pogrubiona bezszeryfowa: 14pt, 13pt, 12pt) [5]
\titleformat{\chapter}[hang]{\normalfont\sffamily\bfseries\fontsize{14}{16}\selectfont}{\thechapter.}{1em}{}
\titleformat{\section}{\normalfont\sffamily\bfseries\fontsize{13}{15}\selectfont}{\thesection}{1em}{}
\titleformat{\subsection}{\normalfont\sffamily\bfseries\fontsize{12}{14}\selectfont}{\thesubsection}{1em}{}

% Konfiguracja bibliografii (BibLaTeX + Biber)
\usepackage[
backend=biber,          % Silnik sortujący (Biber jest nowoczesny i obsługuje UTF-8)
style=numeric,          % Styl cytowania: numeryczny [1], [2] (alternatywa: alphabetic lub authoryear)
sorting=nty,            % Sortowanie w spisie: n-nazwisko, t-tytuł, y-rok
defernumbers=true       % Prawidłowe numerowanie przy podziale na sekcje
]{biblatex}

% Wskazanie pliku bazy danych zawierającego wpisy bibliograficzne
\addbibresource{thesis-references.bib}
